
\subsubsection{Расчет параметров волны в глубокой воде}
\label{sec:deep_whater_wave}

На сегодняшний день существует широкий спектр методов расчета параметров ветрового волнения на глубокой воде -- от эмпирических и параметрических (Pierson–Moskowitz~\cite{pierson_moskowitz_1964}, JONSWAP~\cite{hasselmann_jonswap_1973}, различных региональных fetch-duration кривых~\cite{cerc_spm_1984, shore_fetch_duration_report}) до полноспектральных моделей (WAM~\cite{wamdi_group_1988, ecmwf_wam_doc_2016}, SWAN~\cite{booij_swan_1999}, WAVEWATCH III~\cite{tolman_ww3_2009}) и современных статистических/ML-подходов, требующих расширенного набора входных данных и значительных вычислительных ресурсов. 

Однако в условиях ограниченной исходной информации, представленной лишь полями скорости ветра и длиной разгона волн, наиболее практичным и физически интерпретируемым выбором остается классический метод Свердрупа-Мунка-Бретшнейдера (Sverdrup–Munk–Bretschneider, SMB), основанный на работах~\cite{sverdrup_munk_1947,sverdrup_munk_1947_tagu} и их последующей ревизии~\cite{bretschneider_1958}.
SMB метод напрямую задает зависимость значимой высоты и пикового периода волн от ветра и длины разгона и может быть надежно применен при минимуме доступных данных.

Это эмпирический параметрический метод, калиброванный по большому массиву натурных наблюдений, подробно описанный в Shore Protection Manual (USACE)~\cite{cerc_spm_1984} и современных обзорах~\cite{aisjah2016smb_java_sea, palmieri_2022_from_wind_to_wave}.

В SMB методе характеристики волн (значимая высота волн $H_s$ и период $T_p$) определяются по трем основным параметрам:
\begin{equation}
  H_s = f(U, F, t),
\end{equation}
\begin{equation}
  T_p = g(U, F, t),
\end{equation}
где \(U\) -- скорость ветра на высоте 10 м,
\(F\) -- длина разгона ветра (fetch), \(t\) -- продолжительность ветра.

В основу метода заложена идея, что ветер передает воде энергию вдоль ограниченного разгона и чем длиннее разгон и сильнее ветер, тем выше и длиннее волны, развивающиеся вплоть до предела \enquote{полностью развитого волнения}.
Возникающего, когда подвод энергии уравновешивается диссипацией (рассеиванием энергии путем перехода организованного колебания воды в тепловую и мелкомасштабную турбулентную энергию).

Указанное условие формирует следущее дифференциальное уравнение:
\begin{equation}
    \frac{\partial E}{\partial t} + \frac{\partial}{\partial x}(c_g E) = S_{in} + S_{nl} + S_{ds},\label{eq:base_diff_eq}
\end{equation}
где 
$E(F,t)$ -- средняя волновая энергия на единицу площади, 
$c_g$ -- групповая скорость (скорость переноса энергии),
$S_{in}$ -- приток энергии от ветра,
$S_{nl}$ -- перераспределение энергии нелинейными взаимодействиями волн,
$S_{ds}$ -- диссипативный член, который обычно отрицателен и описывает потери энергии за счет вязкости, турбулентности, обрушения волн и~др.

Для установившегося волнения (стационарный режим, $\partial E / \partial t = 0$) и при фиксированной  длине разгона волны $F$, вместо выражения~\eqref{eq:base_diff_eq} часто рассматривают уравнение баланса в виде:
\begin{equation}\label{eq:diff_balabnc}
    \frac{\partial}{\partial F}(c_g E) = S_{in}(U, E) + S_{ds}(E, \ldots).
\end{equation}

В SMB модель качественно основана на идеи баланса~\eqref{eq:diff_balabnc}, предполагая, что рост волн на разогоне прекращается там, где прирост от ветра $S_{in}$ по величине равен потерям $S_{ds}$, останавливаясь за счет компенсации диссипацией подкачки энергии волны ветром.

%%%%%%%%%%%%%%%%%%%%%%%%%%%%%%%%%%%%%%%%%%%%%%%%

В предпосылках SMB не фигурируют глубина (глубокая вода), начальное волнение и~т.\,п. -- считается, что они либо \enquote{убраны} условиями, либо мало влияют в принятом диапазоне.

В одномерном приближении баланс энергии волн вдоль направления ветра определяется по выражению~\eqref{eq:diff_balabnc}.
 
После перехода от энергии к высоте $E \propto \rho g H_s^2$ и размерностного масштабирования по $U$ и $g$ (ускорение свободного падения), естественное решение в безразмерных переменных принимает вид:
\begin{equation}
    \frac{g H_s}{U^2} = \Phi\!\left(\frac{g F}{U^2}\right),
\end{equation}
где $\Phi$ -- некоторая универсальная функция, определяемая из наблюдений.

Результирующие формулы SMB являются эмпирической аппроксимацией этой зависимости, где по натурным данным была определена кривая в координатах $\bigl(X,\, g H_s/U^2\bigr)$, где:
\begin{equation}
    X = \frac{g F}{U^2}.
\end{equation}

$X$ -- безразмерная величина играющая роль «масштаба разгона» показывая, насколько велика длина разгона волны относительно естественного масштаба, задаваемого ветром и гравитацией.
  
Функция $\Phi$ в SMB описывается формулой вида: 
\begin{equation}
    \frac{g H_s}{U^2} = 0.283\,\tanh\!\bigl(0.53\,X^{0.75}\bigr),
\end{equation}
что в явном виде дает выражение для значимой высоты волн $H_s$ (средняя высота одной трети наибольших волн).

\begin{equation}
    H_s = 0.283\,\frac{U^2}{g}\, \tanh\!\bigl(0.53\,X^{0.75}\bigr),\label{eq:smb_hs}
\end{equation}
где $\tanh(\cdot)$ -- гиперболический тангенс. 

При малых $X$ (для коротких длин разгона ветра) аргумент функции гиперболического тангенса мал, и рост $H_s$ происходит степенным образом с увеличением $F$. 
При больших $X$ функция $\tanh(\cdot)$ стремится к единице, и значимая высота асимптотически приближается к предельному значению:
\begin{equation}
    H_{s,\mathrm{FDS}} \approx 0.283\,\frac{U^2}{g},\label{eq:smb_hs_approx}
\end{equation}
соответствующему полностью развитому ветровому волнению при данном значении скорости ветра $U$.

Аналогично пиковый период волны $T_p$ (с), соответствующий максимуму спектра энергии, выражается через безразмерный $X$ как:
\begin{equation}
    T_p = 7.54\,\frac{U}{g}\, \tanh\!\bigl(0.833\,X^{0.375}\bigr).\label{eq:smb_tp}
\end{equation}

При малых значениях $X$ период невелик и увеличивается с ростом длины разгона $F$, а при больших $X$ также стремится к предельному значению: 
\begin{equation}
    T_{p,\mathrm{FDS}} \approx 7.54\,\frac{U}{g},\label{eq:smb_tp_approx}
\end{equation}
характеризующему полностью развитое ветровое волнение. 

В прибрежных районах, особенно при сложной топографии и неоднородной береговой линии (наличие холмов, городской застройки, мысов и бухт), систематические отличия между береговыми измерениями ветра и ветром над морской поверхностью могут существенно отличаться. 
Это связано с экранирующим эффектом рельефа и различиями в шероховатости подстилающей поверхности. 
Так как участвующие в решении данные представляют скорость ветра, приведенную к высоте 10~м над условной земной поверхностью, необходимо их скорректировать его с учетом того, что скорость над открытой водой выше в силу гладкости водной поверхности.

На практике скорость ветра $U$ в модели может задаваться на основе измерений в км/ч с последующим переводом в м/с и учетом поправочного множителя для условий над водой: 
\begin{equation}
    U_{\text{ms}} = \frac{U_{\text{kmh}}}{3.6}\cdot \text{overwater\_factor},
\end{equation}
где $U_{\text{kmh}}$ -- измеренная скорость ветра в км/ч, 
$\texttt{overwater\_factor}$ -- безразмерный коэффициент, учитывающий отличия между прибрежными измерениями и фактической скоростью ветра над водной поверхностью. 

Руководство Coastal Engineering Manual (USACE)~\cite{usace_cem_2006} рекомендует значение поправки $\texttt{overwater\_factor}\approx1.2$ для метеостанционных данных, тогда как WMO~\cite{wmo_wind_conversion_2010} рассматривает диапазон $1.05–1.20$ в зависимости от стабильности атмосферы. 
В выполняемых расчетах используется среднее консервативное значение $\texttt{overwater\_factor}$ равное 1.1.

Подставляя $U = U_{\text{ms}}$ и заданную длину разгона $F$ в приведенные выше формулы~\eqref{eq:smb_hs} и~\eqref{eq:smb_tp}, получаем значимую высоту $H_s$ и пиковый период $T_p$ ветрового волнения в глубокой воде как функции скорости и длины разгона ветра.

Несмотря на свое распространение, метод SMB не лишен недостатков.
Так в работах \cite{Bishop1992_SPM_WavePredictionReview, hurdle1989revision} указывается, что в общем случае результаты расчета по этому методу дают завышенные значения параметров волн.
Такая ошибка приведет к незначительному завышению нагрузок на берег, что в результате повысит потенциальные значения риска для рассматриваемой территории.
Не смотря на это, в контексте выполняемой работы, такое допущение может быть принято, так как оно формирует некоторую \enquote{страховку} от принятия излишне оптимистичных решений и является меньшим злом относительно риска потенциально неверных выводов из-за случайного занижения риска.  
